\section*{الفصل \ref{ch:g11:lenses-and-eye} --- العدسات والصور والعين}

\begin{solution}{exo:g11:lenses-and-eye:1}
$C = 1/0.50 = +2.0$~D؛ $1/(-0.20) = -5.0$~D؛ $1/0.0040 = +250$~D.
\omterm{def:g11:lenses-and-eye:lens}{وعدسة} قدرها $+8.0$~D: $f' = 1/8.0 = \qty{0.125}{m} = \qty{12.5}{cm}$.
\end{solution}

\begin{solution}{exo:g11:lenses-and-eye:2}
\omterm{def:g11:lenses-and-eye:lens}{مفرِّقة} ($f' < 0$). ولا: فمهما تكن مسافة الجسم $d$،
يقع $\gamma = f'/(f' - d)$ (كما بُيّن في \cref{exo:g11:lenses-and-eye:15})
بين $0$ و$1$ --- \omterm{def:g11:lenses-and-eye:images}{فالصورة
خيالية} معتدلة \emph{مصغَّرة} دائمًا. فيبدو المطبوع معتدلًا
منكمشًا، لا مكبَّرًا أبدًا.
\end{solution}

\begin{solution}{exo:g11:lenses-and-eye:3}
يقع الجسم عند $2f'$؛ ويتقاطع الشعاعان المركزي والموازي
على بعد \qty{20}{cm} خلف \omterm{def:g11:lenses-and-eye:lens}{العدسة}. وللتحقق:
$1/\overline{OA'} = 1/10 - 1/20 = 1/20$، ومنه
$\overline{OA'} = +\qty{20}{cm}$ و$\gamma = 20/(-20) = -1$: أي \omterm{def:g11:lenses-and-eye:images}{حقيقية}
مقلوبة بالحجم الطبيعي --- كما رُسمت.
\end{solution}

\begin{solution}{exo:g11:lenses-and-eye:4}
$1/\overline{OA'} = 1/8.0 - 1/24 = 1/12$:
$\overline{OA'} = +\qty{12}{cm}$، $\gamma = 12/(-24) = -0.50$. أي \omterm{def:g11:lenses-and-eye:images}{صورة
حقيقية} على بعد \qty{12}{cm} خلف \omterm{def:g11:lenses-and-eye:lens}{العدسة}، مقلوبة، بنصف الحجم.
\end{solution}

\begin{solution}{exo:g11:lenses-and-eye:5}
تشكّل \omterm{def:g11:lenses-and-eye:eye}{القرنية} \omterm{def:g11:lenses-and-eye:eye}{والعدسة البلورية} العدسةَ \omterm{def:g11:lenses-and-eye:lens}{المجمِّعة}؛ \omterm{def:g11:lenses-and-eye:eye}{والشبكية} هي
الشاشة. ويغيّر \omterm{def:g11:lenses-and-eye:eye}{التكيّف} \omterm{def:g11:lenses-and-eye:focal}{البعد البؤري} (أي \omterm{def:g11:lenses-and-eye:vergence}{التقارب})
\omterm{def:g11:lenses-and-eye:eye}{للعدسة البلورية}؛ أما المسافة بين \omterm{def:g11:lenses-and-eye:lens}{العدسة} \omterm{def:g11:lenses-and-eye:eye}{والشبكية} فلا تقدر أن تتغير. \omterm{def:g11:lenses-and-eye:eye}{والنقطة
القريبة}: أقرب نقطة تُرى حادّة عند تمام \omterm{def:g11:lenses-and-eye:eye}{التكيّف} ---
وهي \qty{25}{cm} في \omterm{def:g11:lenses-and-eye:eye}{العين المعيارية}.
\end{solution}

\begin{solution}{exo:g11:lenses-and-eye:6}
المنظر: الحسّاس في المستوي البؤري، على بعد \qty{50}{mm} من \omterm{def:g11:lenses-and-eye:lens}{العدسة}. وعند
\qty{3.0}{m}: $1/\overline{OA'} = 1/50 - 1/3000 = 59/3000$، ومنه
$\overline{OA'} \approx \qty{50.8}{mm}$. والانتقال: نحو \qty{0.8}{mm}.
\end{solution}

\begin{solution}{exo:g11:lenses-and-eye:7}
$1/\overline{OA'} = 1/6.0 - 1/5.0 = -1/30$:
$\overline{OA'} = \qty{-30}{cm}$ --- أي \omterm{def:g11:lenses-and-eye:images}{خيالية} معتدلة،
$\gamma = (-30)/(-5.0) = +6.0$: أي ستة أضعاف الطابع، على بعد \qty{30}{cm}
فوق \omterm{def:g11:lenses-and-eye:lens}{العدسة}.
\end{solution}

\begin{solution}{exo:g11:lenses-and-eye:8}
$1/f' = 1/40 + 1/40 = 1/20$: $f' = \qty{20}{cm}$؛
$\gamma = 40/(-40) = -1$. ويفرض التناظر
$\abs{\overline{OA'}} = \abs{\overline{OA}}$، ومنه
$\abs{\gamma} = 1$؛ \omterm{def:g11:lenses-and-eye:images}{والصورة حقيقية} مقلوبة، ومن ثم $\gamma = -1$:
أي بالحجم الطبيعي بالضبط.
\end{solution}

\begin{solution}{exo:g11:lenses-and-eye:9}
$1/\overline{OA'} = -1/25 - 1/50 = -3/50$:
$\overline{OA'} \approx \qty{-16.7}{cm}$،
$\gamma = (-16.7)/(-50) = +0.33$: أي \omterm{def:g11:lenses-and-eye:images}{خيالية} معتدلة بثلث
الحجم. فنظارة قصير النظر \omterm{def:g11:lenses-and-eye:lens}{مفرِّقة}: وكل ما يُرى عبرها
--- بما في ذلك عينا لابسها، من الخارج --- مصغَّر.
\end{solution}

\begin{solution}{exo:g11:lenses-and-eye:10}
صورة اللانهاية عند \omterm{def:g11:lenses-and-eye:eye}{النقطة البعيدة}: $f' = \qty{-0.40}{m}$، ومنه
$C = -2.5$~D. وتقع صورة الجبل على بعد \qty{40}{cm} أمام
\omterm{def:g11:lenses-and-eye:lens}{العدسة} --- أي عند \omterm{def:g11:lenses-and-eye:eye}{النقطة البعيدة} بالضبط، وتراها العين المرتخية حادّة
بلا \omterm{def:g11:lenses-and-eye:eye}{تكيّف}.
\end{solution}

\begin{solution}{exo:g11:lenses-and-eye:11}
$C = 1/\overline{OA'} - 1/\overline{OA} = 1/(-0.80) - 1/(-0.25)
= -1.25 + 4.00 = +2.75$~D.
\end{solution}

\begin{solution}{exo:g11:lenses-and-eye:12}
\emph{1.} يعطي $\gamma = \overline{OA'}/\overline{OA} = -50$
$\overline{OA'} = -50\,\overline{OA}$؛ ومع
$\overline{OA'} - \overline{OA} = \qty{5.1}{m}$،
$-51\,\overline{OA} = \qty{5.1}{m}$: أي $\overline{OA} = \qty{-0.10}{m}$،
$\overline{OA'} = +\qty{5.0}{m}$.

\emph{2.} $1/f' = 1/5.0 + 1/0.10 = 10.2$: $f' \approx \qty{9.8}{cm}$.

\emph{3.} $50 \times \qty{24}{mm} = \qty{1.2}{m}$، مقلوبة --- ولهذا
تُحمَّل الشرائح مقلوبة.
\end{solution}

\begin{solution}{exo:g11:lenses-and-eye:13}
\emph{1.} $C_\infty = 1/0.017 \approx 59$~D؛
$C_{25} = 1/0.017 + 1/0.25 \approx 63$~D. \omterm{def:g10:signals-and-waves:amplitude}{والسعة}:
$1/0.25 = 4.0$~D.

\emph{2.} \omterm{def:g10:signals-and-waves:amplitude}{السعة} $= 1/1.0 = 1.0$~D: أي أن الباقي ربع ($25\%$)
\omterm{def:g10:signals-and-waves:amplitude}{سعة} \omterm{def:g11:lenses-and-eye:eye}{العين المعيارية} البالغة $4.0$~D.
\end{solution}

\begin{solution}{exo:g11:lenses-and-eye:14}
\emph{1.} $f' = 0.25/3 \approx \qty{0.083}{m} = \qty{8.3}{cm}$.

\emph{2.} $\overline{OA'} = \qty{-25}{cm}$:
$1/\overline{OA} = -1/25 - 3/25 = -4/25$، ومنه
$\overline{OA} = \qty{-6.25}{cm}$ و
$\gamma = (-25)/(-6.25) = +4.0$ --- أي أفضل من $\times 3$
المعلَن، الذي يفترض الصورة في اللانهاية؛ فجرّها إلى
\omterm{def:g11:lenses-and-eye:eye}{النقطة القريبة} يكسب \omterm{def:g10:orders-of-magnitude:unit}{وحدة} واحدة.
\end{solution}

\begin{solution}{exo:g11:lenses-and-eye:15}
$1/\overline{OA'} = 1/f' - 1/d = (d - f')/(f'd)$، ومنه
$\overline{OA'} = f'd/(d - f')$ و
$\gamma = \overline{OA'}/(-d) = f'/(f' - d)$. فإذا كان $d > f'$:
$\overline{OA'} > 0$ (أي \omterm{def:g11:lenses-and-eye:images}{حقيقية}) و$\gamma < 0$ (أي مقلوبة). وإذا كان
$0 < d < f'$: $\overline{OA'} < 0$ (أي \omterm{def:g11:lenses-and-eye:images}{خيالية}) و
$\gamma = f'/(f' - d) > 1$: أي معتدلة مكبَّرة --- وهو وضع
\omterm{def:g11:lenses-and-eye:magnifier}{المكبّرة}.
\end{solution}

\begin{solution}{pb:g11:lenses-and-eye:1}
\textbf{1.} $1/\overline{OA'} = 1/20 - 1/60 = 1/30$:
$\overline{OA'} = +\qty{30}{cm}$، $\gamma = -0.50$ --- أي \omterm{def:g11:lenses-and-eye:images}{حقيقية} مقلوبة
بنصف الحجم.

\textbf{2.} $\overline{OA'} = +\qty{60}{cm}$، $\gamma = -2$: أي أن
الموضعين \qty{30}{cm} و\qty{60}{cm} تبادلا الدورين، والتكبيران
مقلوب أحدهما الآخر --- فالضوء المعكوس يتبع \omterm{def:g10:relative-motion:trajectory}{المسار}
نفسه.

\textbf{3.} $1/\overline{OA'} = 1/20 - 1/10 = -1/20$:
$\overline{OA'} = \qty{-20}{cm}$، $\gamma = +2$ --- أي \omterm{def:g11:lenses-and-eye:images}{خيالية} معتدلة
مضاعفة.

\textbf{4.} من اللانهاية تبدأ الصورة في المستوي البؤري؛ وكلما اقترب
الجسم من $F$ تراجعت \omterm{def:g11:lenses-and-eye:images}{الصورة الحقيقية} المقلوبة إلى اللانهاية
وكبرت؛ وداخل $f'$ تنقلب \omterm{def:g11:lenses-and-eye:images}{خيالية} معتدلة مكبَّرة أمام
\omterm{def:g11:lenses-and-eye:lens}{العدسة}.

\textbf{5.} $1/\overline{OA} \approx 0$ يترك
$1/\overline{OA'} = 1/f'$: أي $\overline{OA'} = f'$، وهو المستوي البؤري.

\textbf{6.} في المستوي البؤري، على بعد \qty{50}{mm} خلف \omterm{def:g11:lenses-and-eye:lens}{العدسة}.

\textbf{7.} $1/\overline{OA'} = 1/50 - 1/1000 = 19/1000$:
$\overline{OA'} \approx \qty{52.6}{mm}$؛ والانتقال $\approx \qty{2.6}{mm}$.

\textbf{8.} $1/\overline{OA} = 1/55 - 1/50 = -1/550$: فأقرب مسافة تبؤر
\qty{550}{mm} = \qty{55}{cm}.

\textbf{9.} $\gamma = 55/(-550) = -0.10$: يُصوَّر الوجه على
\qty{20}{mm}، أي داخل الحسّاس \qty{24}{mm} بقليل.

\textbf{10.} يفرض $\gamma = -1$ أن $\overline{OA'} = -\overline{OA}$،
ومنه $2/\overline{OA'} = 1/f'$: أي الجسم على \qty{100}{mm} ($= 2f'$) أمامها،
والحسّاس على \qty{100}{mm} خلفها --- أي ضعف سحب المنظر، وهو أبعد بكثير من
آلية \qty{55}{mm}؛ وعدسات الماكرو توفّر الانتقال الزائد.

\textbf{11.} كل ما هو أقرب من نقطتها القريبة (\qty{1.0}{m}) يتشوّش:
فمدّ الذراعين يقرّب الصفحة إلى \qty{1.0}{m} قدر الإمكان ---
فتكون حادّة وإن صغرت.

\textbf{12.} $C = 1/(-1.0) - 1/(-0.25) = -1.0 + 4.0 = +3.0$~D.

\textbf{13.} $\gamma = (-1.0)/(-0.25) = +4.0$: أي أكبر أربع مرات لكنها
أبعد أربع مرات --- فالحجم الزاوي نفسه؛ والمكسب ليس الحجم بل
\emph{الحدّة}، لأن الصورة صارت الآن عند نقطتها القريبة.

\textbf{14.} تقتضي الحدّة أن تقع الصورة وراء \qty{1.0}{m}: أي من الصفحة
على \qty{25}{cm} (والصورة على \qty{1.0}{m}) إلى الصفحة على
$f' = 1/3.0 \approx \qty{0.33}{m}$ (والصورة في اللانهاية). أما الجسم البعيد
فيُصوَّر على \qty{33}{cm} \emph{خلف} النظارة --- أي \omterm{def:g11:lenses-and-eye:images}{صورة حقيقية}
لا تقدر العين على استعمالها: فيتشوّش العالم البعيد.

\textbf{15.} ومن هنا العدسات ثنائية \omterm{def:g11:lenses-and-eye:focal}{البؤرة}: قدرة القراءة في نصف
\omterm{def:g11:lenses-and-eye:lens}{العدسة} الأسفل، ولا قدرة (أو تصحيح البعد) في نصفها الأعلى.

\textbf{16.} $1/\overline{OA'} = 1/5.0 - 1/4.0 = -1/20$:
$\overline{OA'} = \qty{-20}{cm}$، \omterm{def:g11:lenses-and-eye:images}{خيالية} معتدلة، $\gamma = +5.0$.

\textbf{17.} $1/\overline{OA} = -1/25 - 1/5.0 = -6/25$:
$\overline{OA} \approx \qty{-4.2}{cm}$؛ $\gamma = (-25)/(-25/6) = +6.0$.

\textbf{18.} الحجر عند $F$: فتقع الصورة في اللانهاية، وتُشاهَد بلا
\omterm{def:g11:lenses-and-eye:eye}{تكيّف} --- أي ساعات من الفحص دون إجهاد العين.

\textbf{19.} $C_\infty = 1/0.017 \approx 59$~D؛
$C_{25} = 1/0.017 + 1/0.25 \approx 63$~D: أي نحو $4$~D من
\omterm{def:g11:lenses-and-eye:eye}{التكيّف}.

\textbf{20.} $f' = 1/58.8 \approx \qty{17.0}{mm}$ مقابل
$1/62.8 \approx \qty{15.9}{mm}$: أي نحو \qty{1.1}{mm}، بنسبة $6\%$ ---
وهو التقريب المدمج الذي حاكته آلة التصوير بانتقال قدره \qty{2.6}{mm}
ورقّعته الجدّة بمقدار $+3$~D.
\end{solution}
